\subsection*{调试与详细模式：打印、断言与 GDB}

调试的基本方法是\textbf{逐步缩小问题范围}。先做静态分析和编译时断言，在不增加运行时开销的情况下发现错误；再用有针对性的运行时检查定位问题；范围缩小后，按需使用交互式调试器深入检查。

Croqtile 在每一层都有工具：\textbf{编译时形状打印}（\co@print!@、\co@println!@），不用启动核函数就能核对分块尺寸；\textbf{运行时断言}（\co@assert@），用来检查不变式；\textbf{带守卫的设备 \texttt{println}}，在运行时查看特定线程；给 \co@debugger@ 用的\textbf{调试 RTTI}；以及查看编译器实际执行步骤的\textbf{详细模式}（\co@-v@）。

先使用容易计算参考结果的小输入。矩阵某个元素错误时，还原它的分块坐标和块内坐标，检查选中的输入切片，再比较一个部分和，最后才看完整输出。程序挂起时，则追踪一个事件槽位，记录初始状态、每次触发以及应消费该事件的实例。这样可以区分索引错误与交接缺失。

若失败发生在编译阶段，先确定阶段再修改核函数。预处理输出显示宏展开结果，生成源码显示 Croqtile 的翻译，详细输出则给出目标编译命令。修复后，关闭诊断输出，再次与参考结果比较，然后测量性能。第 1.4 和 2.14 节给出对应命令与检视语法。
